\section*{Hoofdstuk \ref{ch:g12:stretch-reflex} --- De rekreflex en de zenuwboodschap}

\begin{solution}{exo:g12:stretch-reflex:1}
Ontvanger (spierspoeltje), sensorisch \omterm{def:g12:stretch-reflex:neuron}{neuron}, schakelcentrum (grijze stof
van het ruggenmerg, één \omterm{def:g12:stretch-reflex:synapse}{synaps}), motorisch \omterm{def:g12:stretch-reflex:neuron}{neuron}, uitvoerder (de spier).
\end{solution}

\begin{solution}{exo:g12:stretch-reflex:2}
Een cellichaam met de kern, dendrieten die signalen ontvangen, en één axon
(vaak met myeline) dat de boodschap van het \omterm{def:g12:stretch-reflex:neuron}{neuron} naar zijn uiteinden
draagt. De boodschap reist van de dendrieten en het cellichaam langs het
axon naar de uiteinden.
\end{solution}

\begin{solution}{exo:g12:stretch-reflex:3}
Een korte omkering van de membraanspanning, van \qty{-70}{mV} tot ongeveer
\qty{+30}{mV} en binnen een milliseconde terug, die ongewijzigd langs het
axon reist. Alles of niets: onder de drempel niets, erboven altijd
hetzelfde volledige signaal. De sterkte wordt vastgelegd in de frequentie
van de \omterm{prop:g12:stretch-reflex:actionpotential}{actiepotentialen}.
\end{solution}

\begin{solution}{exo:g12:stretch-reflex:4}
\omterm{prop:g12:stretch-reflex:actionpotential}{Actiepotentialen} bereiken het uiteinde; blaasjes geven de \omterm{def:g12:stretch-reflex:synapse}{neurotransmitter}
in de spleet af; die bindt ontvangers op de ontvangende \omterm{def:g10:cells-common-unit:cell}{cel} en verandert
haar spanning (naar de drempel toe of ervan weg); hij wordt daarna
vernietigd of teruggenomen.
\end{solution}

\begin{solution}{exo:g12:stretch-reflex:5}
Acetylcholine. Curare bezet haar ontvangers op de spier zonder hen te
activeren; het botulinegif verhindert haar afgifte uit de blaasjes. Beide
verlammen.
\end{solution}

\begin{solution}{exo:g12:stretch-reflex:6}
6, 12 en 24 pieken in \qty{100}{ms}: 60, 120 en 240 per seconde. De
grootte en de vorm van de pieken veranderen niet.
\end{solution}

\begin{solution}{exo:g12:stretch-reflex:7}
Geleidingstijd $30 - 12 = \qty{18}{ms}$ over \qty{1.5}{m}: ongeveer
\qty{83}{m/s}.
\end{solution}

\begin{solution}{exo:g12:stretch-reflex:8}
De achterwortel draagt de sensorische vezels: doorgesneden bereikt de
boodschap van de rek het merg nooit, maar het motorische \omterm{def:g12:stretch-reflex:neuron}{neuron} en zijn
vezel zijn heel en kunnen nog altijd worden geprikkeld om de spier te
laten samentrekken. De voorwortel draagt de motorische vezels:
doorgesneden kan er helemaal niets meer vanuit het merg bij de spier
komen.
\end{solution}

\begin{solution}{exo:g12:stretch-reflex:9}
Als de tegenwerker ook samentrok, zou hij de beweging tegenwerken en
zichzelf uitrekken, waardoor zijn eigen \omterm{def:g12:stretch-reflex:reflex}{reflex} zou worden uitgelokt --- een
touwtrekwedstrijd. Het tussenneuron van het merg doet het remmen, via een
remmende \omterm{def:g12:stretch-reflex:synapse}{synaps} op de motorische \omterm{def:g12:stretch-reflex:neuron}{neuronen} van de tegenwerker.
\end{solution}

\begin{solution}{exo:g12:stretch-reflex:10}
1 eenheid: niets (onder de drempel). 4 en 8 eenheden: \omterm{prop:g12:stretch-reflex:actionpotential}{actiepotentialen} van
dezelfde \qty{100}{mV}, maar met hogere frequenties, hoger bij 8 dan bij 4.
\end{solution}

\begin{solution}{exo:g12:stretch-reflex:11}
Alleen het uiteinde bevat blaasjes met transmitter en alleen het
ontvangende membraan bevat ontvangers, dus kan de chemische stof maar in
één richting werken. De transmitter afgeven, zijn tocht door de spleet en
de reactie van de ontvangers kosten ongeveer een milliseconde.
\end{solution}

\begin{solution}{exo:g12:stretch-reflex:12}
Acetylcholine hoopt zich op in de verbindingen en laat de vezels blijven
vuren: trekkingen, dan aanhoudende samentrekking en uitputting van de
spieren, ook die van de ademhaling --- de dood door verstikking. Een
middel dat de ontvangers blokkeert (een atropineachtig middel) vermindert
de overmatige prikkeling; te weinig laat de vergiftiging staan, te veel
verlamt zoals curare, dus moet de dosis kloppen.
\end{solution}

\begin{solution}{exo:g12:stretch-reflex:13}
Elke sensorische ingang prikkelt haar motorische \omterm{def:g12:stretch-reflex:neuron}{neuronen} en remt normaal
die van de tegenwerkers via tussenneuronen. Met de remming geblokkeerd
prikkelt een aanraking beide kanten van elk \omterm{def:g10:muscles-and-joints:joint}{gewricht} en verspreidt de
prikkeling zich ongehinderd door het merg: alle spieren trekken tegelijk
samen, en de kramp rekt andere spieren uit waarvan de \omterm{def:g12:stretch-reflex:reflex}{reflexen} er nog bij
komen.
\end{solution}

\begin{solution}{exo:g12:stretch-reflex:14}
Geleidingstijd $60 - 12 = \qty{48}{ms}$ over \qty{1.5}{m}: ongeveer
\qty{31}{m/s}, een derde van het normale. Myeline laat de \omterm{prop:g12:stretch-reflex:actionpotential}{actiepotentiaal}
tussen de openingen in de schede door springen in plaats van
onafgebroken te reizen; zonder haar geleidt de vezel als een dunne vezel
zonder myeline.
\end{solution}

\begin{solution}{exo:g12:stretch-reflex:15}
Het merg draagt boodschappen tussen de hersenen en het lichaam, maar het
verwerkt hen ook: zijn grijze stof bevat de synapsen van de \omterm{def:g12:stretch-reflex:reflex}{reflexbogen},
waar een rek in een samentrekking wordt omgezet en, via tussenneuronen, in
de remming van de tegenwerker, zonder dat er een boodschap bij de hersenen
komt. Een dier waarvan het merg van zijn hersenen is afgesneden, trekt nog
altijd een poot terug bij een kneep en vertoont nog altijd de knieschop:
het merg is een centrum, niet alleen een kabel.
\end{solution}

\begin{solution}{pb:g12:stretch-reflex:1}
\textbf{1.} $0.75 + 0.75 = \qty{1.5}{m}$.

\textbf{2.} $30 - 4 - 4 - 1 = \qty{21}{ms}$ geleiding: $1.5/0.021
\approx \qty{70}{m/s}$.

\textbf{3.} Ongeveer \qty{10}{cm} meer elke kant op, \qty{0.2}{m} in
totaal: $0.2/70 \approx \qty{3}{ms}$ meer, zo'n \qty{33}{ms}.

\textbf{4.} Het elektrische signaal is het op gang brengen van de vezels;
de samentrekking zelf --- het langs elkaar schuiven van de draadjes en de
trek aan de \omterm{def:g10:muscles-and-joints:muscle}{pees} --- heeft enkele tientallen milliseconden nodig om genoeg
kracht op te bouwen om het been te bewegen.

\textbf{5.} Een boodschap van het merg naar de hersenen en een bevel terug
reizen \qty{1.2}{m} met ongeveer \qty{70}{m/s} --- \qty{17}{ms} --- plus
verscheidene synapsen en het verwerken door de hersenen zelf: minstens
\qty{50}{ms}, en een beslissing veel langer. De schop is voorbij voordat
enig signaal van de hersenen de motorische \omterm{def:g12:stretch-reflex:neuron}{neuronen} kan bereiken.

\textbf{6.} Bij 200 per seconde 2 \omterm{prop:g12:stretch-reflex:actionpotential}{actiepotentialen} in \qty{10}{ms} (0.2 in
rust). Het tempo verandert; de grootte en de vorm van elke piek niet.

\textbf{7.} Met een factor 10.

\textbf{8.} Bij 20 per seconde liggen de dosissen transmitter
\qty{50}{ms} uit elkaar en dooft elke dosis vóór de volgende; het
motorische \omterm{def:g12:stretch-reflex:neuron}{neuron} haalt de drempel nooit. Bij 200 per seconde komen zij
\qty{5}{ms} uit elkaar aan en tellen zij op: de drempel wordt gehaald en
het motorische \omterm{def:g12:stretch-reflex:neuron}{neuron} vuurt.

\textbf{9.} Er halen minder motorische \omterm{def:g12:stretch-reflex:neuron}{neuronen} de drempel en elk vuurt
minder \omterm{prop:g12:stretch-reflex:actionpotential}{actiepotentialen}: er trekken minder motorische eenheden samen, met
een lager tempo, en de schop is zwakker.

\textbf{10.} \omterm{prop:g12:stretch-reflex:actionpotential}{Actiepotentialen} zijn alles of niets: hun grootte ligt vast
in het \omterm{def:g12:stretch-reflex:neuron}{neuron}, dus dragen zij geen enkele informatie over de prikkel.
Alleen het tempo kan verschillen.

\textbf{11.} Acetylcholine; nadat hij de ontvangers heeft gebonden, wordt
hij binnen enkele milliseconden door een \omterm{def:g11:enzymes-and-phenotype:enzyme}{enzym} in de spleet vernietigd.

\textbf{12.} De ontvangers van de \omterm{def:g10:muscles-and-joints:muscle}{spiervezel}: de acetylcholine wordt
afgegeven maar kan niet werken.

\textbf{13.} De afgifte van acetylcholine uit de blaasjes. In het
laboratorium zou acetylcholine die rechtstreeks op de spier wordt
aangebracht haar onder botulinegif laten samentrekken (ontvangers vrij)
maar niet onder curare (ontvangers geblokkeerd).

\textbf{14.} De \omterm{prop:g12:stretch-reflex:actionpotential}{actiepotentialen} van de sensorische vezel en, omdat er
niets bij het merg komt, die van het motorische \omterm{def:g12:stretch-reflex:neuron}{neuron} en het
elektromyogram; de spier trekt nog altijd samen als haar zenuw onder de
blokkade wordt geprikkeld.

\textbf{15.} Curare blokkeert het ontvangen, het botulinegif de afgifte,
het verdovingsmiddel de geleiding: drie stappen, één stille spier.

\textbf{16.} De sensorische vezel prikkelt ook een tussenneuron dat de
motorische \omterm{def:g12:stretch-reflex:neuron}{neuronen} van de kniebuigers remt; hun rustactiviteit wordt
onderdrukt, en de tegenwerker ontspant tot onder zijn rustspanning.

\textbf{17.} De \omterm{def:g12:stretch-reflex:reflex}{reflex} heeft alleen het merg nodig, dat op die hoogte heel
is. Zijn overdrevenheid laat zien dat de hersenen normaal signalen sturen
die de \omterm{def:g12:stretch-reflex:reflex}{reflex} dempen; zonder hen draait de boog op volle gevoeligheid.

\textbf{18.} Het spoeltje, de sensorische zenuw of de achterwortel, de
grijze stof van het merg op die hoogte, de motorische zenuw of de
voorwortel, of de spier zelf. De motorische zenuw rechtstreeks prikkelen
en de spier opnemen zegt of de motorische kant werkt; de sensorische zenuw
tijdens een rek opnemen zegt of de sensorische kant werkt.

\textbf{19.} Hij toetst met één tik de hele boog --- sensorische zenuw,
merg, motorische zenuw, spier --- en het dempen ervan door de hersenen:
een ontbrekende, zwakke of overdreven \omterm{def:g12:stretch-reflex:reflex}{reflex} wijst eerder dan enig ander
teken op een niveau van het zenuwstelsel.

\textbf{20.} Ongeveer \qty{70}{m/s}; de frequentie van de
\omterm{prop:g12:stretch-reflex:actionpotential}{actiepotentialen} van het spoeltje; het binden van acetylcholine aan de
ontvangers van de spier.
\end{solution}
